\documentclass{article}
\usepackage{graphicx} % Required for inserting images
\usepackage{amsfonts}

\title{Apostol Chapter 1 - Exercise 1}
\author{Afonso Vilhena Francisco }
\date{06 August 2026}

\begin{document}

\maketitle

\textbf{1. If} $\mathbf{(a,b)=1}$ \textbf{and if} $\mathbf{c |a}$ \textbf{and} $\mathbf{d|b}$, \textbf{then} $\mathbf{(c,d)=1}$

\vspace{0.5cm}

Obviously, $1|c$ and $1|d$, i.e, $1$ is a common divisor of $c$ and $d$.

Let $e \in \mathbb{Z}$ be such that $e|c$ and $e|d$.

Since $e|c$ and $c|a$, by transitivity, we have that $e|a$.
Simmilarly, from $e|d$ and $d|b$, we get that $e|b$. Then, by definition of $(a,b)$, we conclude that $e|(a,b)$, but $(a,b)=1$, then $e|1$.

Therefore, $(c,d)=1$.

\end{document}
